% ============================================================
%  ch08  第三讲 §3.1：Beilinson--Drinfeld Grassmann 簇
%  原文行 1766–2013
% ============================================================

\section{第三讲：Beilinson--Drinfeld Grassmann 簇与分解结构}\label{Lecture III}

本讲引入曲线上的仿射 Grassmann 簇的版本，即 Beilinson--Drinfeld Grassmann 簇，并讨论它们的一些基本几何性质，特别是分解性质。

\subsection{Beilinson--Drinfeld Grassmann 簇}

Beilinson--Drinfeld 的重大发现是：经由定理 \ref{glob aff} 对仿射 Grassmann 簇的解释可以有极为广泛的推广。

首先，可以让定理 \ref{glob aff} 中的点 $x$ 变动。对点 $x\in X(R)$，以 $\Ga_x$ 记映射 $x:\Spec R\to X$ 的图像。于是可定义预叠
\[\Gr_{\underline G, X}(R)=\left\{(x,\mE,\beta)\ \left|\ \begin{array}{l}
x\in X(R),\quad \mE \text{ 是 } X_R \text{ 上的 } \underline{G}\text{-主丛},\\
\beta: \mE|_{X_R-\Ga_x}\cong \mE^0|_{X_R-\Ga_x} \text{ 是平凡化}
\end{array}\right.\right\}.\]
但事实上还可以做得更多：允许有多个点。设 $I$ 是有限非空集。对点 $x=\{x_i,i\in I\}\in X^I(R)$，以
\[\Ga_x:=\bigcup_{i\in I} \Ga_{x_i}\]
记各映射 $x_i: \Spec R\to X$ 的图像之并（作为 $X_R$ 的闭子概形）。

\begin{dfn}
在曲线自积 $X^I$ 上定义预层
\begin{equation}\label{BD Grass def}
\Gr_{\underline G, X^I}(R)=\left\{(x,\mE,\beta)\ \left|\ \begin{array}{l}
x\in X^I(R),\quad \mE \text{ 是 } X_R \text{ 上的 } \underline{G}\text{-主丛},\\
\beta: \mE|_{X_R-\Ga_x}\cong \mE^0|_{X_R-\Ga_x} \text{ 是平凡化}
\end{array}\right.\right\}.
\end{equation}
以 $q_I:\Gr_{\underline G,X^I}\to X^I$ 记结构映射。\zy{$\Gr_{\underline G,X^I}$ 是「把抛物结构放在 $I$ 个点上」的族；\S\ref{Ran space} 用 Ran 空间把这些族按点集之并粘成 $\Gr_{\Ran}$，分解性质即在该对象上表述。nLab 词条 Ran space 与 factorization algebra 记的正是这一框架。[来源：nLab Ran space（2020-07-27）、factorization algebra（2026-08-16）。]}
\end{dfn}

\begin{thm}\label{representability of BD}
预层 $\Gr_{\underline G, X^I}$ 可由 $X^I$ 上有限型的归纳概形表示。若 $\underline G$ 约化，则它归纳射影。
\end{thm}

\begin{proof}
同样可用类似命题 \ref{lin rep} 与 \ref{gen G} 的论证化归到 $\underline G=\GL_n$ 的情形。然后只需证明对正整数 $N$，
\[\Gr_{\GL_n,X^I}^{(N)}=\left\{(x,\mE,\beta)\in\Gr_{\GL_n,X^I}\mid \mO_{X_R}^n(-N\Ga_x)\subset \mE\subset \mO_{X_R}^{n}(N\Ga_x)\right\}\]
可由 $X^I$ 上射影概形表示。用引理 \ref{projectivity} 中的论证，映射 $\mE\mapsto \mF:=\mO_{X_R}^{n}/\mE(-N\Ga_x)$ 把它等同于下述函子：给 $R$ 配上对 $(x,\mF)$，其中 $x\in X(R)$，而 $\mF$ 是 $R$-平坦的 $\mO_{X_R}$-模，且是 $\mO_{X_R}^n/\mO_{X_R}^n(-2N\Ga_x)$ 的商。其可表示性随即由商概形（代替引理 \ref{Embedding} 中的 Grassmann 簇）的可表示性推出。
\end{proof}

在不引起混淆时，把 $\Gr_{\underline G,X^I}$ 简记作 $\Gr_{X^I}$。

\begin{rmk}\label{projectivity of BD}
定理 \ref{representability of BD} 的最后一个论断在更一般的情形下也成立。称 $X$ 上光滑仿射群概形 $\underline G$ 为 \textbf{Bruhat--Tits 群概形}（或\textbf{抛孔群概形}），若它在每个 $D_x$（$x\in |X|$）上的限制都是 Bruhat--Tits 意义下的抛孔群概形。用命题 \ref{indproper}，类似 \cite{PZ} 的论证表明 $\Gr_{X^I}$ 在 $X^I$ 上归纳真。
\end{rmk}

类似地可定义多点版本的射流群与环路群。对 $x\in X^I(R)$，以 $\hat{\Ga}_x$ 记图形 $\Ga_x$ 在 $X_R$ 中的形式完备化（它是形式概形）。有典范映射 $\pi:\hat{\Ga}_x\to X_R$。
\begin{equation}
(L^+\underline G)_{X^I}=\left\{(x,\beta)\mid x\in X^I(R),\ \beta\in
\underline{G}(\hat{\Ga}_x)\right\}.
\end{equation}

\begin{prop}
$(L^+\underline{G})_{X^I}$ 可由 $X^I$ 上仿射概形表示。
\end{prop}

为严格定义整体环路群，需要下述引理。

\begin{lem}
形式概形 $\hat{\Ga}_x$ 归纳仿射。以 $A_x$ 记其函数环的拓扑代数，$\hat{\Ga}'_x=\Spec A_x$。有典范映射 $p:\hat{\Ga}_x\to\hat{\Ga}'_x$，则存在唯一态射 $i: \hat{\Ga}'_x\to X_R$ 使 $\pi=ip:\hat{\Ga}_x\to X_R$。
\end{lem}

\begin{proof}见 \cite[Proposition 2.12.6]{BD}。
\end{proof}

于是 $\hat{\Ga}_x^\circ=\hat{\Ga}'_x\setminus \Ga_x$ 有定义，且
\begin{equation}
(L\underline G)_{X^I}=\left\{(x,\beta)\mid x\in X^I(R),\ \beta\in
\underline G(\hat{\Ga}^\circ_x)\right\}.
\end{equation}

由定理 \ref{descent lemma} 与推论 \ref{global interpretation of loop} 中类似的论证，有

\begin{prop}
作为 $k$-空间，$\Gr_{\underline G,X^I}\cong[(L\underline G)_{X^I}/(L^+\underline G)_{X^I}]$。特别地，
\begin{enumerate}
\item $(L\underline G)_{X^I}$ 可由 $X^I$ 上归纳概形表示；
\item $(L\underline G)_{X^I}$ 自然作用在 $\Gr_{\underline G,X^I}$ 上。
\end{enumerate}
\end{prop}

Beilinson--Drinfeld Grassmann 簇确实是极不寻常的对象。先分析 $I=\{1\}$ 与 $I=\{1,2\}$ 的情形。

设 $I=\{1\}$，并先假定 $\underline G=G\times X$ 常值。对每个 $x\in X$，选取同构 $\mO_x\simeq \mO$，便由定理 \ref{glob aff} 得到典范同构
\[\Gr_{X}|_x=\Gr_{\underline G, x}\cong \Gr_{G}.\]
然而整体上一般没有典范等同 $\Gr_X\neq \Gr_G\times X$。相反，$\Gr_X$ 可看作扭曲积 $X\tilde\times\Gr_G$（扭曲积的定义见 \S\ref{Intro:torsor}）。确切地说，$X$ 上有 $\Aut(D)$-主丛 $\hat{X}\to X$，分类 $x\in X$ 连同一个 $D_x\cong D$：
\[\hat{X}(R)=\left\{(x,\al)\mid x\in X(R),\ \al: \hat{\Ga}'_x\cong \Spec R[[t]]\right\}.\]
同样有 $\Aut^+(D)$-主丛 $\hat{X}^+$，由那些限制为典范同构 $\Ga_x\simeq \Spec R$ 的 $\al$ 组成。回忆 $\Aut(D)$ 作用在 $\Gr_G$ 上（由 \eqref{auto of disc} 定义）。于是有同构
\begin{equation}\label{III: twisted constr}
\Gr_X\cong \hat{X}\times^{\Aut(D)}\Gr_G\cong \hat{X}^+\times^{\Aut^+(D)}\Gr_G.
\end{equation}
第一个同构实际上蕴含 $\Gr_X$ 上存在联络（另见命题 \ref{III:property}）。第二个同构允许把
\begin{equation}\label{III:SV over X}
\Gr_{\leq \mu,X}: =X\tilde\times \Gr_{\leq \mu}:=\hat{X}^+\times^{\Aut^+(D)}\Gr_{\leq \mu}
\end{equation}
定义为闭子概形。等价地，$\Gr_{\leq \mu,X}$ 可等同于 $\Gr_X$ 中分类那些满足 $\inv_x(\beta)\leq \mu$ 的点 $(x,\mE,\beta)$ 的闭子集（赋既约概形结构）。

\begin{rmk}
也可以用上述扭曲手续描述 $(L^+G)_X$ 与 $(LG)_X$。
\end{rmk}

仍设 $I=\{1\}$，但 $\underline G$ 不必常值。例如可以构造 $X$ 上的群概形 $\underline G$，使得 $\underline G|_{X-x}\simeq G\times (X-x)$，而 $\underline G|_{D_x}$ 是 $G$ 的 Iwahori 群概形。此时 $\Gr_{\underline G, X}$ 是从仿射 Grassmann 簇到 $G$ 的仿射旗簇的形变（术语见注 \ref{terminology}）：$\Gr_{\underline G, X}$ 在 $X$ 上点 $y\neq x$ 处的纤维同构于 $\Gr_G$，在 $x$ 处的纤维同构于 $\Fl_G$。这个形变在 Gaitsgory 构造中心层时起关键作用 \cite{Ga}。类似的 $p$-进整数环上的形变在 \cite{PZ} 中被用来定义志村簇的局部模型。

下面转向 $\Gr_{X^2}$。设 $\Delta:X\to X^2$ 是对角。

\begin{prop}\label{III:Gr2}
有典范同构
\[\Delta:\Gr_X\cong \Gr_{X^2}\times_{X^2,\Delta}X,\quad c:\Gr_{X^2}|_{X^2-\Delta}\cong (\Gr_X\times\Gr_X)|_{X^2-\Delta}.\]
此外 $\Gr_{X^2}$ 上有自然 $S_2$-作用，使 $c$ 为 $S_2$-等变，其中 $S_2$ 通过交换两个因子作用在 $\Gr_X\times\Gr_X$ 上。
\end{prop}

\begin{proof}
第一个同构显然。为构造第二个，设 $(x,y)\in (X^2-\Delta)(R)$。则 $(X_R-\Ga_x)$ 与 $(X_R-\Ga_y)$ 构成 $X_R$ 的 Zariski 覆盖。定义映射
\[\Gr_{X^2}\times_{X^2}\Spec R\to \Gr_X\times_{X,x}\Spec R\]
为把 $(\mE,\beta)$ 送到 $(\mE_x,\beta_x)$：其中 $\mE_x$ 由 $\mE|_{X_R-\Ga_y}$ 与 $\mE^0|_{X_R-\Ga_x}$ 经 $\beta$ 粘成，因而带有平凡化 $\beta_x: \mE|_{X_R- \Ga_x}\cong \mE^0|_{X_R-\Ga_x}$。同样有映射 $\Gr_{X^2}\times_{X^2}\Spec R\to \Gr_X\times_{X,y}\Spec R$。合起来定义映射 $\Gr_{X^2}|_{X^2-\Delta}\to (\Gr_X\times\Gr_X)|_{X^2-\Delta}$。

反之，若给定 $(\mE_x,\beta_x)\in \Gr_X\times_{X,x}\Spec R$ 与 $(\mE_y,\beta_y)\in \Gr_X\times_{X,y}\Spec R$，可用 $\beta^{-1}_y\beta_x$ 把 $\mE_{x}|_{X_R- \Ga_y}$ 与 $\mE_y|_{X_R-\Ga_x}$ 粘成 $\mE$，按构造它在 $X_R-\Ga_x-\Ga_y$ 上自然带有平凡化 $\beta$。

最后，非平凡元 $\sigma\in S_2$ 通过把 $(x,y, \mE,\beta)$ 送到 $(y,x,\mE,\beta)$ 作用在 $\Gr_{X^2}$ 上。容易验证在 $c$ 之下，该作用对应于 $\Gr_X\times\Gr_X$ 两个因子的交换。
\end{proof}

这些同构与直觉相悖：$\Gr_{X^2}\to X^2$ 在对角上的纤维比一般纤维「更小」！这是纯粹的无穷维现象，因为我们都知道代数簇族的纤维维数在特殊化时会上升。

事实上并不矛盾，有下述结果。

\begin{prop}\label{III:closure}
设 $\underline G=G\times X$ 常值，$p\nmid \pi_1(G_\der)$。经由上述等同，把 $\Gr_{\leq \mu,X}\times \Gr_{\leq \la,X}|_{X^2-\Delta}$ 看作 $\Gr_{X^2}|_{X^2-\Delta}$ 的子集，并以 $\Gr_{\leq (\mu,\la)}$ 记它在 $\Gr_{X^2}$ 中的闭包。则 $\Gr_{\leq (\mu,\la)}|_{\Delta}= \Gr_{\leq \mu+\la,X}$。
\end{prop}

于是只要限制到有限维子簇，直觉就恢复了。该结果的略弱版本在 \cite{Z07} 中证明（另见 \cite{Z10}），但大致相同的论证在此也适用。

下面考虑线丛。设 $\mL$ 是 $\underline G$-主丛的预叠 $\Bun_{\underline G}$ 上的线丛，且在平凡 $\underline G$-主丛处刚性化。注意有自然遗忘映射
\begin{equation}\label{III:I-un}
u_{X^I}: \Gr_{X^I}\to \Bun_{\underline G}.
\end{equation}
以 $\mL_{X^I}$ 记 $\mL$ 到 $\Gr_{X^I}$ 的拉回。

\begin{prop}\label{III:line}
设 $\underline G=G\times X$ 常值，$G$ 单连通。则在命题 \ref{III:Gr2} 给出的等同下，有典范同构
\[\Delta^*\mL_{X^2}\cong \mL_{X}, \quad \mL_{X^2}|_{X^2-\Delta}\cong c^*(\mL_X\boxtimes\mL_X|_{X^2-\Delta}).\]
\end{prop}

\begin{proof}
第一个同构显然，说明第二个。由命题 \ref{III:Gr2} 的证明，下图交换
\[\begin{tikzcd}
(\Gr_X\times X)|_{X^2-\Delta}\arrow[rr, "c^{-1}(\id\times e)"]\arrow[dr] & & \Gr_{X^2}|_{X^2-\Delta}\arrow[dl] \\
& \Bun_G &
\end{tikzcd}\]
故 $\mL_{X^2}|_{X^2-\Delta}$ 在 $\Gr_X\times X|_{X^2-\Delta}$ 上的限制典范同构于 $\mL_X\boxtimes\mO_{X}|_{X^2-\Delta}$。利用 $\Pic^e(\Gr_G)\simeq \bZ$ 离散（故代数等价于零的线丛都平凡）以及 $\Gr_X$ 连通，可见对 $\Gr_X$ 的每个点 $s$，$\mL_{X^2}|_{X^2-\Delta}$ 在 $\Gr_X\times\{s\}$ 上的限制同构于 $\mL_X$；同理它在 $\{s\}\times \Gr_X$ 上的限制也同构于 $\mL_X$。

由于 $\Gr_X$ 在 $X$ 上归纳真，可用 See-Saw 原理（见 \cite[\S 5]{Mum}）推出 $\mL_{X^2}|_{X^2-\Delta}$ 典范同构于 $\mL_X\boxtimes\mL_X|_{X^2-\Delta}$，且与刚性化相容。
\end{proof}

注意若 $G$ 不单连通，命题 \ref{III:line} 对 $\Bun_G$ 上任意线丛不能成立\footnote{不过，若要求 $\Bun_G$ 上的 $\mL$ 在适当意义下「可乘」，命题大概可以推广到不必单连通的群。}。此外，该论证对非常值群概形 $\underline G$ 也不适用，即使其一般纤维单连通，因为我们关键用到了 $\Gr_{X^I}\to X^I$ 归纳真（以便应用 See-Saw 原理）。多数应用可用下文命题 \ref{det factorization} 代替。

确切地说，\S\ref{determinant line} 中 $\Gr_{\GL_n}$ 上行列式线丛的构造立刻推广，给出 $\Gr_{\GL_n,X^I}$ 上的行列式线丛 $\mL_{\det, X^I}$。

\begin{prop}\label{det factorization}
对行列式线丛 $\mL_{\det,X^I}$，有与命题 \ref{III:line} 类似的同构。
\end{prop}

特别地，给定表示 $\rho:\underline G\to \GL_n$，沿 $f_\rho: \Gr_{\underline G,X^I}\to\Gr_{\GL_n,X^I}$ 拉回，可在 $\Gr_{\underline G,X^I}$ 上构造带命题 \ref{III:line} 那种同构的线丛。

\begin{proof}
第一个同构的存在性显然。第二个同构由典范同构
\begin{equation}\label{det functor}
\det(M\oplus N)\cong \det(M)\otimes_R \det(N)
\end{equation}
推出，其中 $M,N$ 是有限投射 $R$-模。
\end{proof}

\begin{rmk}\label{graded factorization}
然而命题 \ref{III:line} 中的同构 $\mL_{X^2}|_{X^2-\Delta}\cong c^*(\mL_X\boxtimes\mL_X|_{X^2-\Delta})$ 与命题 \ref{det factorization} 中相应同构有本质差别。注意 $\mL_{X^2}$ 与 $\mL_{\det,X^2}$ 都带有自然 $S_2$-作用：命题 \ref{III:line} 的同构是 $S_2$-等变的，而命题 \ref{det factorization} 的同构只差一个符号才 $S_2$-等变。符号出现是因为在 \eqref{det functor} 下，典范同构 $\det(M\oplus N)\cong\det(N\oplus M)$ 对应于典范同构 $\det(M)\otimes\det(N)\cong\det(N)\otimes\det(M)$ 差符号 $(-1)^{\on{rk}M\on{rk}N}$。这反映了：取行列式线是从有限投射 $R$-模范畴 $\on{Proj}_R$（张量积取直和）到\emph{分次}可逆 $R$-模范畴 $\Pic^\bZ_R$ 的对称幺半函子。进一步阐发见 \cite{OZ1}。
\end{rmk}

再列出 Beilinson--Drinfeld Grassmann 簇的两条基本性质。

\begin{prop}\label{III:property}~
\begin{enumerate}
\item $\Gr_{X^I}$ 上存在典范联络，即：设 $x,y:S\to X^I$ 满足 $x|_{S_{red}}=y|_{S_{red}}$，则有典范同构 $\Gr_{X^I}\times_{X^I,x}S=\Gr_{X^I}\times_{X^I,y}S$。
\item Beilinson--Drinfeld Grassmann 簇与 \'etale 基变换交换，即：设 $U\to X$ 是 \'etale 态射，则有典范同构 $\Gr_{\underline G,X^I}\times_{X^I}U^I\cong\Gr_{\underline G_U, U^I}$。
\end{enumerate}
\end{prop}

\begin{proof}
(1) 见 \cite[\S 2.8]{BD}；(2) 见 \cite[Lemma 3.2]{Z10}。
\end{proof}

本小节最后引入 Beilinson--Drinfeld Grassmann 簇的几种变体。首先是 $X^J$ 上的卷积 Grassmann 簇。

设 $\phi:J\twoheadrightarrow \{1,2,\ldots,n\}$ 是有限集的满射，$J_i=\phi^{-1}(i)$。对 $x\in X^J$，以 $x_i\in X^{J_i}$ 记相应分量。定义卷积 Grassmann 簇
\begin{equation}\label{III:conv}
\Gr_{\underline G, \phi}=\left\{(x,(\mE_{i},\beta_i)_{i=1,2,\ldots,n})\ \left|\ \begin{array}{l}
x\in X^J(R),\quad \mE_i \text{ 是 } X_R \text{ 上的 } \underline{G}\text{-主丛},\\
\beta_i: \mE_{i}|_{X_R-\Ga_{x_i}}\cong \mE_{i-1}|_{X_R-\Ga_{x_i}},\ \text{其中 } \mE_0=\mE^0
\end{array}\right.\right\}.
\end{equation}
它有时也记作 $\Gr_{X^{J_1}}\tilde\times\Gr_{X^{J_2}}\tilde\times\cdots\tilde\times\Gr_{X^{J_n}}$，因为有投影
\[\pr_i:\Gr_{X^{J_1}}\tilde\times\Gr_{X^{J_2}}\tilde\times\cdots\tilde\times\Gr_{X^{J_i}}\to \Gr_{X^{J_1}}\tilde\times\Gr_{X^{J_2}}\tilde\times\cdots\tilde\times\Gr_{X^{J_{i-1}}}\]
把 $\Gr_{X^{J_1}}\tilde\times\cdots\tilde\times\Gr_{X^{J_i}}$ 实现为 $\Gr_{X^{J_1}}\tilde\times\cdots\tilde\times\Gr_{X^{J_{i-1}}}$ 与 $\Gr_{X^{J_i}}$ 的扭曲积（定义见 \S\ref{Intro:torsor}）。确切地说，有 $(L^+\underline G)_{X^{J_n}}$-主丛
\begin{equation}\label{torsor E}
\bE\to \Gr_{X^{J_1}}\tilde\times\Gr_{X^{J_2}}\tilde\times\cdots\tilde\times\Gr_{X^{J_{n-1}}}\times X^{J_n},
\end{equation}
分类点 $(x,(\mE_{i},\beta_i)_{i=1,2,\ldots,n-1})\in \Gr_{X^{J_1}}\tilde\times\cdots\tilde\times\Gr_{X^{J_{n-1}}}$、点 $y\in X^{J_n}$，以及 $\mE_{n-1}$ 沿 $\hat{\Ga}_y$ 的平凡化；再用该主丛构造扭曲积。

还有卷积映射
\[m_{\phi}: \Gr_{\underline G, \phi}\to \Gr_{X^J},\]
它把 $(x,(\mE_{i},\beta_i)_{i=1,2,\ldots,n})$ 送到 $(x, \mE_n, \beta_1\cdots\beta_n)$。用类似命题 \ref{III:Gr2} 的论证可得

\begin{lem}\label{generic conv}
$m_\phi$ 在 $X^{(\phi)}$ 上限制为同构。
\end{lem}

其次是 Beilinson--Drinfeld Grassmann 簇的对称化版本。（相对）有效 Cartier 除子的定义见 \cite[Tag01WQ, Tag056P]{St}。回忆给定 $d$，第 $d$ 个对称积 $X^{(d)}=X^d/\!\!/S_d$ 分类 $X$ 上的有效 Cartier 除子，即对每个 $k$-代数 $R$，$X^{(d)}(R)$ 是 $X_R$ 上相对有效 Cartier 除子的集合。于是可如下定义对称化的 Beilinson--Drinfeld Grassmann 簇 $\Gr^{(d)}\to X^{(d)}$：
\begin{equation}\label{sym BD Grass}
\Gr^{(d)}=\left\{(D,\mE,\beta)\ \left|\ \begin{array}{l}
D\subset X_R \text{ 是 } d \text{ 次相对有效 Cartier 除子},\\
\mE \text{ 是 } X_R \text{ 上 } \underline{G}\text{-主丛},\ \beta: \mE|_{X_R-D}\cong \mE^0|_{X_R-D} \text{ 是平凡化}
\end{array}\right.\right\}.
\end{equation}
容易看出自然映射
\[\Gr_{X^d}\to \Gr^{(d)}\times_{X^{(d)}}X^d\]
是同构。给定余权 $\mu$，有闭子簇 $\Gr^{(d)}_{\leq \mu}$，且自然映射
\[\Gr_{\leq (\mu,\ldots,\mu)}\to\Gr^{(d)}_{\leq \mu}\times_{X^{(d)}}X^d\]
是同构。卷积映射诱导 $\Gr_{\leq \mu}\tilde\times\cdots\tilde\times\Gr_{\leq \mu}\to \Gr_{\leq \mu}^{(d)}$。

\begin{ex}\label{momega1BD}
设 $G=\GL_n$，$\omega_1=(1,0,\ldots,0)$。则
\[\Gr^{(m)}_{\leq \omega_1}=\left\{ \mE\subset \mO_{X_R}^n\mid \mO_{X_R}^n/\mE \text{ 是秩 } m \text{ 的投射 } R\text{-模} \right\}.\]
进一步设 $X=\bA^1=\Spec k[t]$。则例 \ref{momega1} 与 \ref{conv v.s. springer} 有整体对应物。首先，公式 \eqref{nil to Gr} 实际定义了
\[\gl_n\to \Gr_{\leq \omega_1}^{(n)}.\]
在整体情形构造 \eqref{sat and springer} 的对应物留作练习。此外，\eqref{sat and springer2} 的整体类比是
\[\begin{CD}
\widetilde{\gl}_n@>>> \Gr_{\omega_1,X}\tilde\times\cdots\tilde\times\Gr_{\omega_1,X}@>>> X^n\\
@VVV@VVV@VVV\\
\gl_n@>>>\Gr^{(n)}_{\leq \omega_1}@>>>X^{(n)},
\end{CD}\]
其中 $\widetilde{\gl}_n\to\gl_n$ 是 Grothendieck--Springer 消解（\cite[\S~1]{Yu}），左侧方块是笛卡尔方块。把 $X^n\cong \bA^n$ 与 $\gl_n$ 的标准 Cartan 子代数 $\frakt_n$ 等同，把 $X^{(n)}$ 与 $\frakt_n/\!\!/S_n$ 等同，则下排两个映射的复合可等同于把矩阵送到其特征多项式的 Chevalley 映射。这个图最早出现在 \cite{MVy}。
\end{ex}

\begin{rmk}
$p$-进群的 Beilinson--Drinfeld 仿射 Grassmann 簇版本也存在，最近由 P. Scholze 构造 \cite{Sch}，但已超出本讲义范围。
\end{rmk}
